
\section{Preliminaries}
\sectionframe{2}{Preliminaries: before diving in... some definitions and preliminary results}

%--------------------------
\myframe{Random signals and stochastic processes}{%

\begin{block}{Random signal}    
    A random signal can defined as the tuple $(T,V,S)$ where
    \begin{itemize}
        \item $T$ is the set of time. We consider discrete-time signals, $T=\Z$
        \item $V$ is the \tq{alphabet} of signal values, $V=\R$ or $V=\R^n$
        \item $S$ is the sample space. Every $s\in S$ is a function $s:t\mapsto s(t)$
        where $t\in T$ and $s(t) \in V$ also called \textit{trajectory} of the signal
    \end{itemize}
\end{block}

\pause

\begin{block}{Stochastic process}
    Mathematical formulation/description used to formally model the aleatoric
    nature of a random signal. Consisting of the tuple $(S,\S, P)$ where:
    \begin{itemize}
        \item $S$ sample space. Set of all possible trajectories
        \item $\S$ events set of all logical propositions or events regarding the signal
        \item $P$ probability measure on $\S$, i.e. $P: \S\mapsto [0,1]$ mapping
        every $A\in\S$ to $P(A)\in[0,1]$, the a-priori
        \tq{confidence} $A$ actually happening
    \end{itemize}
\end{block}
}


%--------------------------
\myframe{A more general formulation}{%

\begin{block}{Stochastic process (usual def in prob. theory)}
    A collection of random variables, $\{y(t)\}$, indexed by a temporal index
    $t\in T$, all defined on the same (abstract) probability space $(\Omega, \A,
    P)$ where:
    \begin{itemize}
        \item $\Omega$ event set of \tq{state of nature} $\omega$
        \item $\A$ $\sigma$-algebra of subsets of $\Omega$ 
        \item P probability measure on $\A$
    \end{itemize}
    Each variable is a ($\A$-measurable) function $y(t): t \mapsto y(t,\omega)$. The
    \textit{trajectories} are the successive determinations
    $\{y(\cdot,\tilde{\omega})\}$ of $y(t)$ corresponding to a particular event
    $\tilde{\omega}\in\Omega$
\end{block}

\pause

Our definition $(S,\S,P)$ is a \tq{concrete} version of the abstract
$(\Omega,\A,P)$. If we define a variable
$y(t)\in(S,\S,P)$ as
%
$$
y(t,s) := s(t)
$$
where $y(t)$ associates at every trajectory $s\in S$ its value at $t$, then it
is clear that, for varying $t$, the successive determinations of the variables
$\{y(t)\}$ corresponding to the event $\overline{s}\in \S$ describe indeed the
trajectory $\{\overline{s}(t)\}$ of the process.
}

%--------------------------
\myframe{Let's clarify...}{%

The probabilistic formulation and formalization introduced is \textit{just} a
convenient choice to describe (or model) the unpredictable nature of a signal.
It is worth remembering that this is one possible choice that eventually
allows one to resort to tools from probability theory to describe, handle and
develop useful tools for filtering and identification. For this reason, it is
important to understand that, regardless the real nature of the underline
signal, the same formulation (as \textit{stochastic process}) can indeed be used
to describe every signal (aleatoric or not). On the other hand, it might be
worth remembering that this choice is dictated by the \tq{simplicity},
compactness and generality to describe events that could indeed also be
described using a \tq{deterministic} formulation.
}

%--------------------------
\myframe{The Problem of \textit{Estimation}: Estimate vs Measure}{%

Statistical problem of estimation seen as formalization of the mathematical
problem of measuring where the goal is that of getting information about a
desired quantity $x$, non directly observable, based on observable data, $y$,
whose values are produced by the measurement \tq{instrument}.

\vspace{3mm}

\pause
Ideally, knowing the \tq{physical coupling}, i.e., relation, between $x$ and
$y$ should allow to retrieve the value of $x$ based on observable values of
$y$. Practically, this relation is either not perfectly known or influenced
by \textit{environmental conditions} which introduce \textit{uncertainty} in
the measurement process.

$$
y=h(x) \Longleftrightarrow y=h(x) + \epsilon
$$

As a consequence, the objective can be understood (or interpreted) as that
of \textit{estimating} the value of $x$ based on the observed values of $y$.

\vspace{3mm}

\pause
Following this \tq{uncertainty-based} interpretation:
\begin{itemize}
    \item it is not always straightforward how to generally get $x$ from $y$
    \item being impractical (meaningless) to reconstruct $x$ with certainty, how
    to \textit{best} approximate it? 
\end{itemize}
}

%--------------------------
\myframe{The Bayesian approach}{%

\begin{block}{}
    $x$ assumed to be a random quantity characterized by some a-priori
    information
\end{block}

\pause

\begin{alertblock}{Notation: random variables vs sample values}
    It is important to understand the difference between a random variable, for
    which bold letters $\x$ will be used vs the corresponding sample values,
    i.e., the actual value taken by the random variable, for which normal case
    letters $x$ will be used.
\end{alertblock}

\pause

\begin{block}{Bayesian inference}
    Assume there exist a random variable/vector $\x$ described by the
    probability density $p_\x(x)$ and the random variable/vector $\y$ of
    measurements described by the conditional density $f(y|x) = f(y|\x=x) =
    \frac{P(y\leq \y \leq y+dy | \x=x)}{dy}$.
    
    \vspace{3mm}

    The problem of Bayesian inference/estimation is that of reconstructing $\x$
    based on $p_\x$, $f(y|x)$ and the observation of $\y=y$. This can be
    understood as to:
    \begin{itemize}
        \item compute the new distribution of $\x$ determined by the observation $\y=y$, or
        \item reconstruct $x=\x(\omega)$ determined by the \tq{measurement}
        process
    \end{itemize}
\end{block}
}

%--------------------------
\myframe{Bayesian inference: complete reconstruction}{%
Assuming $p_\x(x)$ and $f(y|x)$ completely known, the problem is that of
computing

$$
f(x|y) = \frac{f(y|x)p_\x(x)}{\int f(y|x)p_\x(x) dx} = \frac{p_{\y\x}(y, x)}{p_\y(y)}
$$

i.e., the conditional \textit{a-posteriori} density of $\x$ given $\y=y$, which
is the most exhaustive and complete description possible.

\vspace{5mm}

\pause

Explicit analytical computation of the above equations often not possible.

\vspace{5mm}

\pause

In practice, this is rarely the case and/or even necessary because rather than
an exhaustive description, a more handy point estimate of $\x$, telling us what
is the unobservable sample value $\x=x$ that gave rise to the observation
$\y=y$, is more useful.

}

%--------------------------
\myframe{Bayesian inference: point estimate}{%

Assume the relation between $\x$ and $\y$ could analytically be described as

$$
\y = h(\x, \w)
$$

being $\w$ the random vector capturing the variable interations with the
environment, i.e., the \tq{measurement noise}. 

\vspace{5mm}

\pause

During the measurement process the particular experimental conditions $\omega$
give raise to a specific determination $\w(\omega)=w$ which, in turn, give raise
to an observed value $\y=y$ linked to $\x=x$ by the relation

$$
y=h(x,w)
$$

\pause

Then, the problem is that of reconstructing the value $x$ of $\x$ that
\tq{caused} $\y=y$ during the measurement process. Since estimating/knowing $w$
is intrinsically impossible, the problem is understood as that of \tq{best}
approximating $x$.

\vspace{3mm}

\pause

It is then necessary to define a \tq{reasonable} class of approximation criteria.
}

%--------------------------
\myframe{Bayesian inference: conditional/posterior expected risk}{%

\begin{block}{Cost/Loss Function}
    It is called \tq{loss/cost} a scalar function $c: \R^n\mapsto\R$ of $\xi=z-x$,
    $z\in\R^n$, s.t. $c\geq0$, $c''>0$ (strictly convex) and $c(0)=0$. It is said
    symmetric if $c(-\xi)=c(\xi)$.

    \vspace{3mm}

    E.g.: $c(z-x)=\|z-x\|^2_Q$ where $Q>0$ symmetric and $\|x\|^2_Q = x^TQx$.
\end{block}

\pause
$z$ plays the role of approximation of $x$. The problem could then be that of
finding $z$ that minimizes $c$ for every value of the unknown $x$. In practice
it is \tq{smarter} to consider all the available probabilistic information,
e.g., given by $\y=y$.

\pause
\begin{block}{Conditional Expected Risk and Point Bayesian Estimate}
    Consider the \textit{Conditional Expected Risk}
    $$
    R(z,y):=\E[c(z-\x)|\y=y]= \int_{\R^n} c(z-x)f(x|y)dx
    $$
    Then, the \textit{Point Bayesian Estimate} is defined as
    $$
    \hat{x}= \textnormal{Arg min}_z R(z,y)
    $$
\end{block}
}

%--------------------------
\myframe{Bayesian inference: Conditional Expectation}{%

\begin{exampleblock}{Theorem: Conditional Expectation}
    If the loss $c$ is symmetric and the conditional density $f(\cdot|y)$ is
    symmetric around its mean $\mu(y)=\E[\x|\y=y]=\int xf(x|y)dx$, then
    $\hat{x}$ does not depend on $c$ and it is equal to the conditional mean of
    $\x$ given $\y=y$, i.e., $\E[\x|\y=y]$.

    \vspace{3mm}

    Moreover, if the cost $c$ is of quadratic form, i.e., $\|\cdot\|^2_Q$, then
    $\hat{x}=\E[\x|\y=y]$ for all types of conditional density $f(x|y)$
    (assuming the conditional mean exists).

    \vspace{3mm}

    More generally, among all possible measurable functions $g:
    \R^m\mapsto\R^n$, for every symmetric $Q>0\in \R^{n\times n}$, 
    % the conditional mean $\E[\x|\y]$ is that minimizing the
    % mean of the squared distance $\|\x-g(\y)\|_Q$, i.e.,
    $$
    \E[\x|\y] = \textnormal{Arg min}_{g(\cdot)} \E[\|\x-g(\y)\|_Q^2]
    $$
\end{exampleblock}

\vspace{1mm}

\pause

Interpreted as a function of $y$, for symmetric $c$ and $f$, the conditional
mean $\E[\x|\y]: y \mapsto \E[\x|\y=y]$ is the Bayesian
estimator\footnote{Estimator: function (algorithm) of the observed data mapping
them to point estimates.} $\hat{x}$ of $x$ that minimizes the conditional
expected risk $R(z,y)$. This is also true for any $f$ if $c$ is quadratic. In
this case it is referred to as the \textit{minimum squared Bayesian estimator}. 
}

%--------------------------
\myframe{Theorem: Conditional Expectation (proof of last part)}{%

For the law of total expectation, we can write
%
\begin{align*}
    \E[\|\x-g(\y)\|_Q^2] &= \int_\R p_\y(y) \E[\|\x-g(\y)\|_Q^2 | \y = y]dy \\
                         &= \E\left[
                                \E[ \|\x-g(\y)\|_Q^2 | \y ]
                               \right] \\
                         &= \E\left[
                                \E[ \|\x - \E[\x|\y] + \E[\x|\y] - g(\y)\|_Q^2 | \y ]
                              \right] \\
                         &= \E\Big[
                                   \E[ [\x - \E[\x|\y]]^T    Q [\x - \E[\x|\y]]    | \y ] \\
                         &\quad + 2\E[ [\x - \E[\x|\y]]^T    Q [\E[\x|\y] - g(\y)] | \y] \\
                         &\quad +  \E[ [\E[\x|\y] - g(\y)]^T Q [\E[\x|\y] - g(\y)] | \y ]
                               \Big]\\
                         &= \E[\|\x-\E[\x|\y]\|_Q^2] + \E[\|\E[\x|\y]-g(\y)\|_Q^2]  
\end{align*}

where the last equality holds by taking the mean and since the second double
product term is zero conditioning on $\y$. And since the two terms are
non-negative and the first does not depend on $g$, we can minimize the second
by setting $g(\y)=\E[\x|\y]$.

}

%--------------------------
\myframe{Bayesian inferece: MAP estimate}{%

\textit{What if, instead of minimizing risk, one wants to maximize gain?}

\vspace{2mm}

\pause
\begin{block}{Expected gain}
    $$
    G(z,y):=\E[(z-\x)|\y=y] = \int_{\R^n}\gamma(z-x)f(x|y)dx
    $$
    being $\gamma$ a concave function with its max in the origin.
\end{block}

\vspace{2mm}

\pause
By choosing $\gamma$ to arbitrarily approximate Dirac's impulse $\delta$
$\longrightarrow$ $G(z,y)\approx f(z|y)$ 

\vspace{2mm}

The corresponding Max Expected Gain (or \textit{Maximum a Posterior} - M.A.P.)
estimator is given by

$$
\hat{x}_{MAP}(y) := \textnormal{Arg max}_z f(z|y)
$$

which, for general $f(x|y)$ is the \textbf{conditional mode} ($\equiv \E[\x|\y]$
for unimodal symmetric densities).

\pause
\begin{alertblock}{}
    Unfortunately analytical computations of mean (or mode) are usually not possible.
\end{alertblock}

}

%--------------------------
\myframe{Bayesian inference: the Gaussian case}{%
% One particularly important case where the computations can be carried out is
% when $\x$ and $\y$ are jointly Gaussian.

\begin{exampleblock}{The Gaussian case}
    If $\x$ and $\y$ are jointly Gaussian described by a density function with
    mean and variance
    $$
    \mu_\z = \begin{bmatrix} \mu_\x \\ \mu_\y \end{bmatrix}; \quad
    \Sigma_\z = \begin{bmatrix} \Sigma_\x & \Sigma_{\x\y}\\ \Sigma_{\y\x} & \Sigma_\y \end{bmatrix}
    $$
    the conditional density is still Gaussian with mean and variance (if
    $\Sigma_\y>0$\footnote{If this is not the case it can be shown that using
    the pseudo-inverse $(\cdot)^\dagger$ works.})
    \begin{align*}
        \E[\x|\y] &= \mu_\x + \Sigma_{\x\y}\Sigma_\y^{-1}(\y-\mu_\y)\\
        \Var[\x|\y] &= \Sigma_\x - \Sigma_{\x\y}\Sigma_\y^{-1}\Sigma_{\y\x}
    \end{align*}
\end{exampleblock}

\pause

Interpreting $\E[\x|\y]$ as the Bayesian estimator of $\x$ given $\y$, notice
that:
\begin{itemize}
    \item the estimation error is defined as: $\tilde{\x}:=\x-\E[\x|\y]$
    \item $\tilde{\x}$ is independent from the data, i.e., all available info
    has been exploited
    \item $\Var[\x|\y]$ (cond. var. of $\x|\y$) can be understood as
    $\Sigma_{\tilde{\x}}$ (uncond. var. of $\tilde{\x}$), does not depend on the
    data, and is given by the difference between $\Sigma_\x$ and
    $\Var\E[\x|\y]=\Sigma_{\x\y}\Sigma_\y^{-1}\Sigma_{\y\x}$ (uncond. vars. of
    $\x,\ \E[\x|\y]$, resp.)
\end{itemize}
}

%--------------------------
\myframe{The Gaussian case: proof}{%

\begin{exampleblock}{}
    Consider $\overline{\z}:=\z-\muz \sim \N(0,\Sigma_{\z})$.
    Consider a linear transformation of $\tilde{\x} = \overline{\x} + A\overline{\y}$,
    $\tilde{\y} = \overline{\y}$ such that $\tilde{\x}\perp\tilde{\y}$. That is
    $$
    \E[\tilde{\x}\tilde{\y}^T] = 0 = \Sigma_{\x\y} + A\Sigma_{\y} \Longrightarrow A=-\Sigma_{\x\y}\Sigma_{\y}^{-1}
    $$
    Clearly, since linear transforms maintain Gaussianity,
    $$
    \tilde{\z}\sim\N(0,\Sigma_{\tilde{\z}})=
        \begin{bmatrix}
            I & A \\ & I
        \end{bmatrix}
        \Sigma_{\z}
        \begin{bmatrix}
            I & \\
            A^T & I
        \end{bmatrix}
        =
        \begin{bmatrix}
            \Sigma_{\tilde{x}} & \\ & \Sigma_{\y}
        \end{bmatrix},\quad
        \Sigma_{\tilde{x}}=\Sigma_{\x} - \Sigma_{\x\y}\Sigma_{\y}^{-1}\Sigma_{\y\x}
    $$
    %
    and being Gaussian and uncorrelated, $\tilde{\x}$ and $\tilde{\y}$ are also
    independent
    \begin{align*}
        \Longrightarrow \overline{\x}=\tilde{\x}-A\overline{\y}
            \Longrightarrow &\E[\overline{\x}|\overline{\y}]=-A\overline{\y}=\Sigma_{\x\y}\Sigma_{\y}^{-1}\overline{\y}\\
                            & \Var[\overline{x}|\overline{\y}]=\Sigma_{\tilde{\x}} = \Sigma_{\x} - \Sigma_{\x\y}\Sigma_{\y}^{-1}\Sigma_{\y\x}
    \end{align*}

    Introducing the means you get the final result.
\end{exampleblock}

}

%--------------------------
\myframe{Conditional Expectation: brief recap}{%

    So far we said that:

    \medskip
    \begin{itemize}
        \pause
        \item $\E[\x|\y]$ is the function $g(\cdot)$ of the data $\y$ that
        minimizes $\E\|\x-g(\y)\|_Q^2$ among the (vast) class of measurable
        functions that are also unbiased, i.e., $\E[g(\y)]=\E[\x]$;
        
        \pause
        \item for $Q=I$, $\E\|\x-g(\y)\|^2_Q$ is the scalar variance of the
        estimation error defined as $\tilde{\x}:=\x-\E[\x|\y]$ $\Longrightarrow$
        $\E[\x|\y]$ is the \textit{minimum variance estimator} of $\x$;
        
        \pause
        \item if $\x$ and $\y$ are jointly Gaussian, the computation of
        $\E[\x|\y]$ (and $\Var[\x|\y]$) can be explicitly performed and depends
        only on first and second moments of the joint density;
        
        \pause
        \item in general, the explicit computation of $\E[\x|\y]$ is
        impractical/impossible, being the conditional expectation a general
        \textit{non-linear} function of the data.
    \end{itemize}

    \vspace{5mm}

    \pause
    What if the functional family of $g(\cdot)$ is a-priori restricted? E.g. to be linear
}

%--------------------------
\myframe{Linear Minimum Mean Squared Error (MMSE) Estimators}{%
Restricting $g(\cdot)$ to the family of linear unbiased functions of $\y$, we
have

\begin{block}{Linear MMSE Estimator: the Problem (vector form)}
    The Linear MMSE Estimator is the linear function of the data
    $$
    g(\y) = A\y + b,\quad A\in\R^{n\times m},\ b\in\R^n
    $$
    that minimizes the mean of the squared estimation error $\E\|\x-g(\y)\|^2$,
    that is
    $$ 
    \min_{A,b} \E\|A\y+b-\x\|^2.
    $$
\end{block}
}

%--------------------------
\myframe{Minimum introduction to Hilbert spaces}{%

A Hilbert space $\H$ is a vector space (over $\R$ or $\C$) equipped with an
inner product $\inner{\cdot}{\cdot}$ which is complete w.r.t. the norm
$\|x\|=\sqrt{\inner{x}{x}}$ induced by the inner product. That is, every Cauchy
sequence converges with respect to $\|\cdot\|$. 

\vspace{5mm}

Examples:
\begin{itemize}
    \item $\R^n$ endowed with the standard Euclidean inner product;
    \item $L^2(\Omega, \A, P)$ of second order variables of a random experiment
    $\{\Omega, \A, P\}$: the space of variables $f$, $\E[|f|^2]<\infty$ (mean
    squared integrable $\equiv$ finite variance), with inner product
    $\inner{\xi}{\eta}=\E[\xi\overline{\eta}]$ given by the correlation between
    variables;
    \item $\H(\y) := \spn{\y_1,\ldots,\y_m} = \left\{\sum_{i=1}^m a_i\y_i;\
    a_i\in\R\right\} \subseteq L^2(\Omega, \A, P)$ \textit{linearly} generated
    by the second order random process $\y$ (collection of variables) with
    $\E[\y_i]=0$ (zero mean) and $\|\y_i\|^2_{\H} = \E[|\y_i|^2]  = \Var[\y_i] <
    \infty$ (finite variance).
\end{itemize}

\vspace{5mm}

\pause
\begin{alertblock}{}
    $\Longrightarrow$ resorting to Hilbert spaces and, in particular, $\H(\y)$ allows to interpret
the problem geometrically!
\end{alertblock}

}


%--------------------------
\myframe{Geometric formulation}{%
    \begin{exampleblock}{Theorem: Orthogonal Projections (scalar form)}
        Let $\x\in\H$ be a (scalar) random variable. Let $\H(\y)\subset\H$. The
        variable $\z^*\in\H(\y)$ minimizing $\|\x-\z\|^2_{\H}=\E[|\x-\z|^2]$ is unique and is
        the orthogonal projection of $\x$ onto $\H(\y)$. Necessary and
        sufficient condition for $\z^*$ to be the orthogonal projection is that
        $$
            \x-\z^* \bot\ \H(\y)\ \ \equiv\ \ \E[(\x-\z^*)\y_i]\ \ i\in\until{m}
        $$
        being $\y_i$ the generators of $\H(\y)$.
    \end{exampleblock}

    \vspace{5mm}

    \pause
    Thanks to this fundamental theorem, the properties inherited from Hilbert
    spaces and the correspondence between the notion of orthogonality ($\bot$)
    and uncorrelation (through the inner product/norm) a solution to the
    \textit{Linear MMSE Estimator} can be drawn (directly in vector form from a
    element-wise application of the projection theorem).

}

%--------------------------
\myframe{Linear Minimum Mean Squared Error (MMSE) Estimators}{%
Restricting $g(\cdot)$ to the family of linear unbiased functions of $\y$, we
have

\pause
\begin{block}{Linear MMSE Estimator: the Problem (vector form)}
    The Linear MMSE Estimator is the linear function of the data
    $$
    g(\y) = A\y + b,\quad A\in\R^{n\times m},\ b\in\R^n
    $$
    that minimizes the mean of the squared estimation error $\E\|\x-g(\y)\|^2$,
    that is
    $$ 
    \min_{A,b} \E\|A\y+b-\x\|^2.
    $$
\end{block}

\pause
\begin{exampleblock}{Linear MMSE Estimator: the Solution (vector form)}
    It can be shown that the solution $\hat{\x}(\y)=g^*(\y)$ (also denoted
    $\Elin[\x|\y]$) is given by
    $$
    \hat{\x}(\y) = \mu_\x + \Sigma_{\x\y}\Sigma_\y^{-1}(\y-\mu_\y)
    $$
    that is $A^*=\Sigma_{\x\y}\Sigma_\y^{-1}$, $b^*=\mu_\x-A^*\mu_\y$.
    Equivalently to the Gaussian case, it depends only of the first
    and second moments of the joint density.    
\end{exampleblock}

}

%--------------------------
\myframe{Notation disclaimer}{%
\customblock{alertblock}{}{$\Elin[\x|\y] \neq \E[\x|y]$}

\vspace{.5cm}

\begin{itemize}
    \item $\E[\x|y]$ minimum variance estimator = conditional mean
    \item $\Elin[\x|\y]$ linear minimum variance estimator $\neq$ conditional mean
    \item $\Elin[\x|\y]$ in general depends only on up to second order statistics
    \item $\Elin[\x|\y]$ sometimes referred to as \textbf{weak conditional mean}
    \item $\Elin[\x|\y] = \E[\x|y]$ for jointly Gaussian densities
    \item $\Elin[\x|\y]$ very \tq{far} from conditional mean for general densities
\end{itemize}

}

%--------------------------
\myframe{Example with code}{%
Before looking at more theory ... a first example of Kalman Filtering!

\vspace{1.5cm}

Jupyter notebook: \texttt{linear\_constant\_velocity.ipynb}

}